% BEGIN REV09_NIVELES_2_5
\subsection{Compatibility of the level-two and level-five modular readers}
\label{corpus9:iv:niveles}

On the modular realization subsequent to the HMT construction,
fix \(\operatorname{Im}\tau>0\), \(q=e^{2\pi i\tau}\) and
\(q^{1/5}=e^{2\pi i\tau/5}\). The Rogers--Ramanujan reader is
\[
 R(q)=\frac{q^{1/5}}{1+\dfrac{q}{1+\dfrac{q^2}{1+\cdots}}},
 \qquad x=R(q)^5.
\]
The modular identity of this realization, with the normalization
of \(j\) used in \eqref{km:j-nivel2}, is
\begin{equation}
 j=-\frac{(x^4-228x^3+494x^2+228x+1)^3}
          {x(x^2+11x-1)^5}.
 \label{corpus9:iv:nivel5-identidad}
\end{equation}
The identity is the modular-realization premise of the
following algebraic transport; the seeds, orientation and
self-scaling representations remain those of APP--TRIT--TPK.

\begin{proposition}[Pentadic quotient and branch selection]
\label{corpus9:iv:nivel5-cociente}
For \(x\ne0\), let \(u=x-x^{-1}\). On the domain \(u\ne-11\),
\begin{equation}
 j=-\frac{(u^2-228u+496)^3}{(u+11)^5}.
 \label{corpus9:iv:nivel5-u}
\end{equation}
The involution \(x\mapsto-1/x\) preserves \(u\). Its preimages are
\[
 x=\frac{u\pm\sqrt{u^2+4}}2.
\]
The quadratic covering ramifies at \(u=\pm2i\); the
subsequent rational map \(u\mapsto j\) has degree six.
\end{proposition}
\begin{proof}
The exact factorizations are
\[
 \begin{aligned}
 x^4-228x^3+494x^2+228x+1
   &=x^2(u^2-228u+496),\\
 x^2+11x-1&=x(u+11).
 \end{aligned}
\]
Substitution into \eqref{corpus9:iv:nivel5-identidad} proves
\eqref{corpus9:iv:nivel5-u}. The fibre equation is
\(x^2-ux-1=0\); its discriminant gives the ramification points.
The numerator of \eqref{corpus9:iv:nivel5-u} has degree six
and the denominator degree five. At \(u=-11\), the
numerator polynomial has value \(3125\), so the two have
no common factor. This proves the asserted degree.
\end{proof}

For real \(u\), the two roots have opposite signs. On the
complex domain, the branch selected in the parametrization
by \(\tau\) is retained. The reading \(j\) alone retains a
fibre with generically six values of \(u\).

\begin{proposition}[Self-scaling limit and compatibility of levels]
\label{corpus9:iv:niveles-ramas}
In the real limit \(q\to1^-\),
\[
 R(q)\longrightarrow\varphi_{\mathrm{HMT}}^{-1},\qquad
 x\longrightarrow\varphi_{\mathrm{HMT}}^{-5}
     =\frac{5\sqrt5-11}{2},\qquad u\longrightarrow-11.
\]
When \(t=t_2(\tau)\) and \(u\) are read on the same modular
parameter, they satisfy
\begin{equation}
 (t+256)^3(u+11)^5
   +t^2(u^2-228u+496)^3=0.
 \label{corpus9:iv:niveles-enlace}
\end{equation}
Writing \(\epsilon=u+11\), the branches satisfying the stated
limits obey
\begin{align}
 \epsilon\to0,\ t\to0:
 &\quad \frac{t^2}{(-\epsilon)^5}\longrightarrow
          \frac{2^{24}}{5^{15}},\\
 \epsilon\to0,\ t\to\infty:
 &\quad t\epsilon^5\longrightarrow-5^{15}.
 \label{corpus9:iv:niveles-asintotica}
\end{align}
\end{proposition}
\begin{proof}
For \(0<q<1\), the alternating positive convergents of the
continued fraction bracket \(R(q)\). As \(q\to1^-\), each
fixed convergent tends to the corresponding convergent of
the continued fraction with all numerators equal to one.
The two subsequences of these convergents have limit
\(\varphi_{\mathrm{HMT}}^{-1}\), in its already established
quadratic realization. The bracketing proves the first limit.
Taking the fifth power and using \(\varphi^2=\varphi+1\)
gives the other two.

Equating \eqref{km:j-nivel2} and \eqref{corpus9:iv:nivel5-u}
proves \eqref{corpus9:iv:niveles-enlace}. Substitution of
\(u=\epsilon-11\) gives
\[
 \frac{(t+256)^3}{t^2}
   =-\frac{(3125-250\epsilon+\epsilon^2)^3}{\epsilon^5}.
\]
On the branch \(t\to0\), the identity rearranges as
\[
 \frac{t^2}{(-\epsilon)^5}
   =\frac{(t+256)^3}{(3125-250\epsilon+\epsilon^2)^3}.
\]
Its limit is \(256^3/3125^3=2^{24}/5^{15}\).
On the branch \(t\to\infty\), multiplication by \(\epsilon^5\)
and expression of the left-hand side as
\(t\epsilon^5(1+256/t)^3\) give the second limit.
\end{proof}

The relation retains the branches of a correspondence over the
same \(j\). The choice of \(\tau\), or of the corresponding
lift, fixes the pair being transported. The pentadic pole
and Fricke inversion retain their respective domains and
operations.
% END REV09_NIVELES_2_5
